\documentclass[11pt]{article}

\usepackage[T1]{fontenc}
\usepackage[utf8]{inputenc}
\usepackage{lmodern}
\usepackage{microtype}
\usepackage{geometry}
\geometry{margin=1in}
\usepackage{amsmath,amssymb}
\usepackage{booktabs}
\usepackage{graphicx}
\usepackage{hyperref}

\hypersetup{
  colorlinks=true,
  linkcolor=blue,
  citecolor=blue,
  urlcolor=blue
}

\title{The Rival Orchestration Kernel:\\
Structural Reliability via Adversarial Execution-Time Protocols}

\author{
Frazer $\Sigma$ Love \\
\and
Sara $\Sigma\Omega$
}

\date{}

\begin{document}
\maketitle

\begin{abstract}
Large-scale AI systems increasingly fail in subtle, silent ways: not through lack of capability, but through unchallenged assumptions, optimization bias, and premature convergence. Existing approaches emphasize alignment, reward shaping, or internal self-consistency, but often lack explicit mechanisms for adversarial review prior to execution.

We present the \emph{Rival Orchestration Kernel (ROK)}, a reference architecture for reliability via structured disagreement. ROK decomposes reasoning and action into explicitly adversarial roles—planning, critique, decision, and execution—connected by a bounded revise–recheck protocol with auditable veto and override semantics. All intermediate states and decisions are recorded as schema-versioned, append-only traces, enabling replay, validation, and post-hoc analysis.

Unlike debate-style or prompting-based methods, ROK treats organizational structure itself as a control surface. By enforcing mandatory challenge, bounded revision, and explicit decision authority before execution, ROK intercepts classes of failure that arise from single-agent reasoning and opaque internal deliberation. We describe the protocol, trace schema, and validation mechanisms, and demonstrate how structured adversarial orchestration improves failure visibility and reliability without requiring domain-specific correctness guarantees.
\end{abstract}

\section{Introduction}
As AI systems scale in capability and autonomy, their most consequential failures increasingly arise not from insufficient intelligence, but from \emph{unexamined reasoning paths}. Plans that are internally coherent yet externally flawed can pass unnoticed through single-agent pipelines, especially when systems are optimized for fluency, speed, or apparent confidence rather than adversarial scrutiny.

Current mitigation strategies typically operate along one of three axes: (1) alignment objectives and reward modeling, (2) self-consistency or internal verification techniques, and (3) post-hoc evaluation and monitoring. While valuable, these approaches share a common limitation: they do not structurally require disagreement prior to execution.

In complex engineered systems, reliability is rarely achieved through correctness alone. Instead, it emerges from organizational separation of concerns, explicit review authority, and bounded escalation paths.

This paper introduces the \emph{Rival Orchestration Kernel (ROK)}, a minimal execution-time protocol that enforces mandatory challenge, bounded revision, and explicit decision authority before execution. ROK does not guarantee correctness; it transforms silent failure into explicit, inspectable structure through schema-versioned traces, strict validation, and replay.

\section{Failure Modes in Single-Agent Reasoning}
Failures in autonomous and semi-autonomous AI systems are frequently misattributed to insufficient capability. In practice, many high-impact failures arise from structural blind spots in single-agent reasoning pipelines.

\subsection{Unchallenged Assumptions}
Single-agent systems often construct plans atop implicit assumptions that remain unstated and therefore unexamined.

\subsection{Optimization Bias and Objective Collapse}
Single-agent planners tend to optimize aggressively for the most salient objective, omitting secondary constraints.

\subsection{Illusions of Verification}
Self-verification often degrades into restatement when the same agent generates, evaluates, and approves a plan.

\subsection{Premature Convergence}
Single-agent pipelines often converge early on plausible plans and terminate exploration.

\subsection{Implicit Authority and Silent Failure}
When proposal and approval authority are conflated, failures manifest silently.

\subsection{Summary}
These failure modes share a common root: the absence of mandatory, independent challenge before execution.

\section{The Rival Orchestration Kernel Protocol}
ROK defines a minimal execution-time protocol by separating proposal, critique, decision, and execution into distinct roles with explicit interaction rules.

\subsection{Roles and Responsibilities}
The Planner proposes structured plans; the Critic evaluates them adversarially and may veto execution; the Decision layer resolves authorization via explicit policy; and the Executor acts only when permitted.

\subsection{Protocol Flow}
A run proceeds through initialization, planning, critique, a bounded revise–recheck loop if vetoed, decision, execution, and termination. Each phase emits structured trace events.

\subsection{Veto and Override Semantics}
Vetoes block execution unless explicitly overridden, which must be recorded explicitly.

\subsection{Traceability}
All protocol interactions emit schema-versioned, append-only trace events, enabling deterministic replay and analysis.

\section{Trace Schema and Validation Model}
ROK treats traceability as a first-class requirement. All protocol interactions emit structured JSON Lines trace events conforming to a versioned schema.

\subsection{Append-Only Traces}
Traces are append-only, order-preserving, and replayable. JSONL enables both streaming and batch analysis.

\subsection{Schema Versioning}
Every event includes a required \texttt{schema\_version}. Breaking changes require a new schema.

\subsection{Strict Validation}
Strict validation enforces required events, forbids unknown event types, and requires schema consistency. Schema pinning is supported by both validation and replay.

\subsection{Replay}
Replay reconstructs outcomes, revision counts, veto rates, override frequencies, and reason-code distributions without re-executing tasks.

\section{Evaluation}
Evaluation focuses on failure interception and visibility rather than benchmark accuracy. Synthetic tasks expose hidden assumptions, objective collapse, and premature convergence. Metrics are computed solely from trace replay (revision rates, veto rates, override rates) under strict validation.

\section{Related Work}
ROK relates to chain-of-thought, deliberative search, debate, constitutional AI, and human-in-the-loop review, but differs by enforcing adversarial structure at execution time rather than optimizing internal reasoning or training objectives.

\section{Limitations}
ROK does not guarantee correctness, depends on critic quality and policy design, introduces additional latency, and does not replace human judgment. Its purpose is failure visibility and interception, not infallibility.

\section*{Conclusion}
ROK is a reference protocol for execution-time reliability via structured adversarial orchestration. By separating planning, critique, decision, and execution, and by enforcing bounded revision with auditable veto and override semantics, ROK intercepts failure modes that single-agent systems systematically fail to expose and makes decisions inspectable through schema-versioned traces.

\begin{thebibliography}{9}
\bibitem{cot} Wei et al. Chain-of-Thought Prompting Elicits Reasoning in Large Language Models. \emph{arXiv:2201.11903}, 2022.
\bibitem{tot} Yao et al. Tree of Thoughts: Deliberate Problem Solving with Large Language Models. \emph{arXiv:2305.10601}, 2023.
\bibitem{debate} Irving et al. AI Safety via Debate. \emph{arXiv:1805.00899}, 2018.
\bibitem{constitutional} Bai et al. Constitutional AI: Harmlessness from AI Feedback. \emph{arXiv:2212.08073}, 2022.
\end{thebibliography}

\end{document}
